\chapter{README }
\hypertarget{md__r_e_a_d_m_e}{}\label{md__r_e_a_d_m_e}\index{README@{README}}
{\bfseries{Huffman Compression Engine}}

This project is a complete C++ implementation of the Huffman Coding algorithm. It provides a full pipeline for compressing and decompressing text files by analyzing character frequency and generating an optimal prefix-\/based binary tree.

The system supports both encryption (compression) and decryption (reconstruction) using a persistent key file, ensuring data can be recovered across different program executions.

{\bfseries{Core Features}}


\begin{DoxyItemize}
\item {\bfseries{Frequency Analysis}}
\item Dynamically calculates the occurrence of each character in the input file to determine optimal encoding weights.
\item {\bfseries{Huffman Tree Construction}}
\item Builds a binary tree using a priority-\/based list, where characters with lower frequency are placed deeper, and frequent characters get shorter codes.
\item {\bfseries{Full Formatting Support}}
\item Preserves spaces, special characters, and newline formatting, ensuring the decrypted output matches the original file exactly
\item {\bfseries{Persistent Key Storage}}
\item Stores the generated Huffman codes in a key file, allowing reconstruction of the tree without needing the original input.
\item {\bfseries{Tree Reconstruction}}
\item Rebuilds the Huffman tree from the saved key file using pointer-\/based traversal logic.
\item {\bfseries{Memory Management}}
\item Implements recursive deletion of tree nodes to prevent memory leaks.
\end{DoxyItemize}

{\bfseries{How It Works}}

{\bfseries{1. Input Processing}}

The program reads the entire text file, including newline characters.

{\bfseries{2. Frequency Calculation}}

Each character is counted and stored in a custom list structure.

{\bfseries{3. Tree Construction}}

The two least frequent nodes are repeatedly combined to form a parent node until a single root remains.

{\bfseries{4. Encoding}}

Each character is assigned a unique binary code based on its path in the tree:\+

Left → 0

Right → 1

{\bfseries{5. Key Generation}}

The character-\/to-\/code mapping is saved in a file (Key\+\_\+for\+\_\+encryption.\+txt).

{\bfseries{6. Encryption}}

The original text is converted into a string of 0s and 1s and stored in Encrypted\+\_\+file.\+txt.

{\bfseries{7. Reconstruction}}

The tree is rebuilt from the key file using the stored codes.

{\bfseries{8. Decryption}}

The encoded bitstream is traversed through the reconstructed tree to recover the original text, saved in Decrypted\+\_\+file.\+txt.

{\bfseries{Technical Implementation}}


\begin{DoxyItemize}
\item Language
\item C++
\end{DoxyItemize}

{\bfseries{Core Data Structures}}


\begin{DoxyItemize}
\item Custom tree node (encryption)
\item Inherited STL list as a priority structure
\item Pointer-\/based binary tree
\end{DoxyItemize}

{\bfseries{Key Techniques}}


\begin{DoxyEnumerate}
\item Recursion (tree traversal, deletion, encoding)
\item File handling using ifstream and ofstream
\item Operator overloading (for frequency increment)
\item Custom sorting for priority handling
\end{DoxyEnumerate}

{\bfseries{Project Workflow}}


\begin{DoxyEnumerate}
\item Read and encode original file
\item Build Huffman tree
\item Save key file
\item Rebuild tree from key
\item Save encrypted file
\item Decrypt file using rebuilt tree
\end{DoxyEnumerate}

{\bfseries{Current Limitations}}


\begin{DoxyItemize}
\item Encoded data is stored as strings ("{}0"{} and "{}1"{}) instead of actual bits
\item This does not reduce physical file size significantly
\item Uses list instead of a more efficient priority\+\_\+queue
\end{DoxyItemize}

{\bfseries{Future Improvements}}


\begin{DoxyEnumerate}
\item {\bfseries{Bit Packing}} Convert binary strings into actual bits for real compression and reduced file size.
\item {\bfseries{Performance Optimization}} Replace list-\/based priority handling with a priority\+\_\+queue (min-\/heap).
\item {\bfseries{User Interface}} Add a command-\/line menu for file selection and operation control.
\item {\bfseries{Error Handling}}
\end{DoxyEnumerate}

\&\+\#x20; Improve robustness for missing files and invalid inputs. 